% !TeX root = slides.tex

% ==========================================================
% USER SETTINGS
% ==========================================================
% Font scheme. Pick ONE. Each scheme pairs a body face, a display face
% (title / frame titles / dividers) and a MATCHING math font, so formulas
% stop looking like they were pasted in from another document.
%
%   helvetica-garamond  Helvetica body, EB Garamond display, Palatino math
%                       (the template's original look)
%   sans                Helvetica throughout + true sans math (newtxsf)
%   garamond            EB Garamond throughout + matching math
%   libertinus          Libertinus throughout + matching math
%
\def\slidefontscheme{helvetica-garamond}

% Letterspacing/protrusion refinements (microtype). Set to 0 to disable.
\def\slidemicrotype{1}

\usepackage{pgfpages}
\setbeamertemplate{note page}{\pagecolor{yellow!5}
%\insertslideintonotes{0.25} \\
\vspace{2em}
\insertnote}
%\setbeameroption{show notes on second screen=right} % Both
\setbeameroption{hide notes} % Only slide
%\setbeameroption{show only notes}
\newcommand{\mathleft}{\@fleqntrue\@mathmargin0pt}
\newcommand{\mathcenter}{\@fleqnfalse}


\usepackage{ulem}
\renewcommand{\ULdepth}{1.8pt}

% PACKAGES ==============================================
\usepackage{booktabs}
\usepackage{pifont}   % \cmark for table indicator cells
\usepackage{multirow}
\usepackage{etoolbox}
\usepackage{type1cm}  % scalable CM sizes: silences "font shape not available"
\usepackage{bbm}
\usepackage{adjustbox}
\usepackage{bigstrut}
\usepackage{soul}
\usepackage{comment}
\usepackage{mathtools}

% TIKZ PROGRAMMING ======================================
\usepackage{tikz}
\usetikzlibrary{matrix}
\usetikzlibrary{shapes,arrows,fit,tikzmark}
\usetikzlibrary{shapes.misc}
\usetikzlibrary{backgrounds}
\usetikzlibrary{shapes.multipart}
\usetikzlibrary{arrows.meta}


%% DOCUMENT ENCODING ====================================================================
\usepackage[utf8]{inputenc} % set input encoding (not needed with XeLaTeX)
\usepackage[T1]{fontenc}

% BIBLIGRAPHY ===========================================
\usepackage{csquotes}
\usepackage[natbib = true, backend = bibtex, style  = authoryear, maxcitenames = 2]{biblatex}
\addbibresource{references.bib}

% IGNORE WARNINGS =======================================
\usepackage{silence}
\WarningFilter{biblatex}{Patching footnotes failed}
\WarningFilter{biblatex}{Using fall-back BibTeX(8) backend:}

% DEFINE FONT ==============================================
% Every scheme must define \displayfamily (used by the title page, frame
% titles and section dividers) and load exactly one math font package.
% Selected via \slidefontscheme at the top of this file.

\newcommand{\displayfamily}{}

% -- Helvetica body + EB Garamond display + Palatino math (original look) --
\ifdefstring{\slidefontscheme}{helvetica-garamond}{%
	\usepackage[scaled=.93]{helvet}
	\renewcommand{\familydefault}{\sfdefault}
	\usepackage[lining]{ebgaramond}
	\usepackage[upint]{newpxmath}
	\renewcommand{\displayfamily}{\fontfamily{EBGaramond-TLF}\selectfont}
}{}

% -- All Helvetica, with genuine sans math so equations match the body --
\ifdefstring{\slidefontscheme}{sans}{%
	\usepackage[scaled=.93]{helvet}
	\renewcommand{\familydefault}{\sfdefault}
	\usepackage[cmintegrals]{newtxsf}
	\renewcommand{\displayfamily}{\sffamily}
}{}

% -- EB Garamond throughout, including matching math --
% \usefonttheme{serif} is required: beamer's default font theme forces
% \sfdefault, which would leave serif math sitting on a sans body.
\ifdefstring{\slidefontscheme}{garamond}{%
	\usefonttheme{serif}
	\usepackage[lining]{ebgaramond}
	\usepackage{ebgaramond-maths}
	\renewcommand{\displayfamily}{\rmfamily}
}{}

% -- Libertinus throughout, including matching math --
\ifdefstring{\slidefontscheme}{libertinus}{%
	\usefonttheme{serif}
	\usepackage{libertinust1math}
	\usepackage[mono=false]{libertinus}
	\renewcommand{\displayfamily}{\rmfamily}
}{}

\ifdefempty{\displayfamily}{%
	\PackageError{preamble}{Unknown \string\slidefontscheme\space
		'\slidefontscheme'. Valid: helvetica-garamond, sans, garamond,
		libertinus}{}%
}{}

\usefonttheme{professionalfonts}  % keep beamer from overriding math fonts

% DOUBLE SPACING ===========================================
\usepackage{setspace}
\setstretch{1.2}

\geometry{
	paperwidth =\the\paperwidth,
	paperheight=\the\paperheight}%   vmargin=0cm,%   head=0.5cm, %   headsep=0pt,%   foot=0.5cm % }

% COLORS ===================================================
% Ratios below are WCAG 2.1 contrast against the white slide background,
% measured, not estimated. Thresholds: >=4.5 body text, >=3.0 large text
% (>=18pt bold, i.e. frame titles) and graphical objects such as plot lines.
%
% -- Semantic roles: use these instead of raw color names -------------
\definecolor{ink}{HTML}{1A1A1A}       % 17.40  body text
\definecolor{inkmuted}{HTML}{595959}  %  7.00  secondary text
\definecolor{inkfaint}{HTML}{6E6E6E}  %  5.10  captions, footline, citations
\definecolor{dimmed}{HTML}{8C8C8C}    %  3.36  de-emphasized LARGE text only
\definecolor{rule}{HTML}{D9D9D9}      %  1.41  DECORATIVE ONLY: hairlines
\definecolor{wash}{HTML}{F5F5F5}      %  1.09  FILL ONLY: panel backgrounds
\definecolor{accent}{HTML}{D10F2F}    %  5.50  primary emphasis
\definecolor{accentalt}{HTML}{004488} %  9.62  secondary emphasis
\definecolor{highlight}{HTML}{FFF0BF} % FILL ONLY: ink on it reads at 15.31

% -- Categorical series for plots and marks --------------------------
% Chosen by exhaustive search: every color clears 3.0 on white, and the set
% stays distinguishable under normal vision, deuteranopia, protanopia and
% tritanopia (worst-case pairwise Lab dE = 23.9). Use in numerical order;
% the ordering is optimized so that any prefix is as distinct as possible.
% Past ~4 series, pair color with direct labels or markers as well.
\definecolor{series1}{HTML}{332288}   % 12.17  indigo
\definecolor{series2}{HTML}{CC6677}   %  3.66  rose   (graphics/large text)
\definecolor{series3}{HTML}{882255}   %  8.73  wine
\definecolor{series4}{HTML}{999933}   %  3.02  olive  (graphics/large text)
\definecolor{series5}{HTML}{4477AA}   %  4.70  blue
\definecolor{series6}{HTML}{543005}   % 11.65  sepia
% Matched light tints for area fills / bars; ink on any of them reads >=13.
\definecolor{fill1}{HTML}{E0DEED}
\definecolor{fill2}{HTML}{F7E8EB}
\definecolor{fill3}{HTML}{EDDEE6}
\definecolor{fill4}{HTML}{F0F0E0}
\definecolor{fill5}{HTML}{E3EBF2}
\definecolor{fill6}{HTML}{E5E0DA}

% -- Legacy names kept so existing decks keep compiling ---------------
% Note: `red`, `orange`, `cyan` and `magenta` shadow xcolor's built-ins, so
% \textcolor{red} here is NOT xcolor's red. Prefer the semantic names above.
% Ratios are on white; those marked (!) are too light for body text.
\definecolor{beamer@blendedblue}{RGB}{25,70,114}  %  9.72
\definecolor{red}{RGB}{209,15,47}                 %  5.50  = accent
\definecolor{lightgray}{HTML}{F5F5F5}             %  1.09  fill only
\definecolor{lightgrey}{RGB}{218,218,218}         %  1.40  fill only
\definecolor{red_slides}{RGB}{187,12,27}          %  6.60
\definecolor{blugreen}{RGB}{0,158,115}            %  3.42  (!) large/graphics
\definecolor{orange}{RGB}{219,107,0}              %  3.43  (!) large/graphics
\definecolor{red_tol}{HTML}{D92120}               %  5.01
\definecolor{blue_tol}{RGB}{0,119,187}            %  4.82
\definecolor{ash}{RGB}{100,100,100}               %  5.92
\definecolor{lightash}{RGB}{140,140,140}          %  3.36  (!) large/graphics
\definecolor{magenta}{RGB}{238,51,119}            %  3.91  (!) large/graphics
\definecolor{cyan}{RGB}{51,187,238}               %  2.21  (!) fill only
\definecolor{green_tol}{RGB}{27,120,55}           %  5.54
\definecolor{purple_tol}{RGB}{118,42,131}         %  8.64
\definecolor{orange_tol}{RGB}{232,96,28}          %  3.43  (!) large/graphics

% TITLE PAGE =============================================
\setbeamertemplate{title page}[default]
\setbeamercolor{title}{fg=beamer@blendedblue,bg=}
\setbeamercolor{author}{fg=black,bg=}
\setbeamercolor{institute}{fg=black,bg=white}
\setbeamercolor{date}{fg=black,bg=}
%\setbeamerfont*{title}{series=\bfseries\selectfont,size=\LARGE}
\setbeamerfont*{title}{size=\huge, series=\bfseries\selectfont, family=\displayfamily}
\setbeamerfont*{institute}{size=\small}

% FRAME TITLE ==============================================
\setbeamercolor{background canvas}{bg=white} 				% background of slides
\setbeamercolor{normal text}{fg=ink}			         	% color of text
%\setbeamercolor{frametitle}{bg=white,fg=blue_slides} 		% slide title color
%\setbeamerfont*{frametitle}{size=\large}
\setbeamerfont{frametitle}{family=\displayfamily, series=\bfseries, size=\LARGE}

% FRAME MARGINS ============================================
\setbeamersize{
	text margin left  = 4ex,
	text margin right = 4ex}

% APPENDIX ===========
\usepackage{changepage}
\usepackage{appendixnumberbeamer}
\renewcommand\appendixname{Appendix} % otherwise weird warning

% FRAME FOOTLINE ===========================================
%\setbeamertemplate{footline}[frame number]
\setbeamertemplate{footline}{
	\leavevmode%
	\hbox{%
	\begin{beamercolorbox}[wd=\paperwidth,ht=2.25ex,dp=1.5em,right]{title in head/foot}%
		\textcolor{inkfaint}{
		\presenter \, $\Big\vert$ \,  \hyperlink{slide:outline}{\insertframenumber\enskip/\enskip\inserttotalframenumber}}  \hspace*{2ex}
%		%\usebeamerfont{date in head/foot} \insertauthor\inserttitle \insertframenumber{} / \inserttotalframenumber \hspace*{2ex}
	\end{beamercolorbox}}%
	\vskip0pt%
}

% ITEMIZE BULLETS ==========================================
\setbeamertemplate{navigation symbols}{}
\setbeamertemplate{itemize items}{$\cdot$}
\setbeamertemplate{itemize subitem}{-}
\setbeamercolor{itemize item}{fg=inkfaint}
\setbeamercolor{itemize subitem}{fg=inkfaint}
\setbeamercolor{enumerate item}{fg=inkfaint}
\setbeamercolor{enumerate subitem}{fg=inkfaint}
\setbeamercolor{button}{bg=white,fg=inkfaint}

\setlength{\leftmargini}{1ex}
\setlength{\leftmarginii}{2ex}
\setlength{\leftmarginiii}{1ex}

% TABLE OF CONTENTS ========================================
\setbeamertemplate{section in toc}{\vspace{2ex}\huge{\alert{\textbf{\textcolor{accent}{\inserttocsectionnumber\,\big|\,\inserttocsection}}}}\vspace*{0.25ex}}
\setbeamertemplate{section in toc shaded}{\enskip\inserttocsection}
\setbeamertemplate{subsection in toc}{\hspace*{-4ex}\enskip\enskip\Large{\inserttocsectionnumber.\inserttocsubsectionnumber\enskip\inserttocsubsection}\par\vspace*{0.5ex}}
\setbeamertemplate{subsection in toc shaded}{\hspace*{-4ex}\enskip\enskip\large{\inserttocsectionnumber.\inserttocsubsectionnumber}\enskip\inserttocsubsection\par}
\setbeamercolor{section in toc shaded}{fg=dimmed}
\setbeamercolor{subsection in toc shaded}{fg=dimmed}

% ALERT COLOR =============================================
\setbeamercolor{alerted text}{fg=accent}

% USER WRITTEN COMMAND ABBREVIATIONS ======================
\newcommand{\alertbf}[1]{\textcolor{accent}{\textbf{#1}}}
\newcommand{\redsmallbold}[1]{{\small\textcolor{accent}{\textbf{\colorbox{white}{#1}}}}}
\newcommand{\citebib}[1]{{\scriptsize\textcolor{inkfaint}{#1}}}
\newcommand{\cornerlinks}[1]{
	\begin{tikzpicture}[remember picture, overlay]
	\node[yshift=0.45ex,anchor=south west]() at (current page.south west){%
		#1 %
	};
	\end{tikzpicture}
}
\newcommand{\cornerinfo}[1]{
	\begin{tikzpicture}[remember picture, overlay]
	\node[yshift=0.35ex,anchor=south west]() at (current page.south west){%
		\textcolor{inkfaint}{\tiny{#1}} %
	};
	\end{tikzpicture}
}
\newcommand{\returnbutton}[1]{%
	\begin{tikzpicture}[remember picture, overlay]
	\node[shift={(-2em,-1.5em)}]() at (current page.north east){%
	\hyperlink{#1}{\beamerreturnbutton{Back}}};
	\end{tikzpicture}%
}

\newcommand{\cmark}{\ding{51}}

\newcolumntype{C}[1]{>{\centering\arraybackslash}p{#1}}
\newcolumntype{L}[1]{>{\raggedright\arraybackslash}p{#1}}
\newcolumntype{R}[1]{>{\raggedleft\arraybackslash}p{#1}}
\newcolumntype{H}{>{\setbox0=\hbox\bgroup}c<{\egroup}@{}}

% MODIFIED FOR ROWCOLOR TO WORK WITH TRANSITIONS
\makeatletter
\def\rowcolor{\noalign{\ifnum0=`}\fi\bmr@rowcolor}
\newcommand<>{\bmr@rowcolor}{%
    \alt#1%
        {\global\let\CT@do@color\CT@@do@color\@ifnextchar[\CT@rowa\CT@rowb}% 
        {\ifnum0=`{\fi}\@gooble@rowcolor}% 
}

\newcommand{\@gooble@rowcolor}[2][]{\@gooble@rowcolor@}
\newcommand{\@gooble@rowcolor@}[1][]{\@gooble@rowcolor@@}
\newcommand{\@gooble@rowcolor@@}[1][]{\ignorespaces}
\makeatother

% PACKAGES TO LOAD LAST ====================================
\usepackage{bm}
\usepackage{pgfplots}
\pgfplotsset{compat = newest}
\usetikzlibrary{positioning}
\usepgfplotslibrary{fillbetween}
\usepackage{amsmath}

% POSITIONING GRID =========================================
% Drop \showgrid into a frame while placing overlay annotations; delete it
% when done. Origin is the slide center, units are cm (16:9 slide = 16x9 cm).
\newcommand{\showgrid}{%
	\begin{tikzpicture}[remember picture, overlay, shift={(current page.center)},
		gl/.style={font=\tiny\bfseries, text=black, fill=white, inner sep=0.6pt}]
		\draw[step=0.5cm, black!10, very thin] (-8,-4.5) grid (8,4.5);
		\draw[step=1cm,   black!10, very thin] (-8,-4.5) grid (8,4.5);
		\draw[black, line width=0.5pt] (-8,0) -- (8,0) (0,-4.5) -- (0,4.5);
		\foreach \x in {-8,...,7} \node[gl] at (\x,0) {\x};
		\foreach \y in {-4,...,4} {\ifnum\y=0\else
			\node[gl] at (0,\y) {\y};\fi}
	\end{tikzpicture}%
}

% CALLOUT BOX ==============================================
% A "key takeaway" panel: soft background, accent rule down the left edge.
\newcommand{\takeaway}[2][0.92\textwidth]{%
	\begin{center}%
	\begin{tikzpicture}
		\node[fill = wash, inner sep = 1.4ex, text width = #1,
		      align = left, text = ink] (tw) {#2};
		\draw[accent, line width = 2.2pt]
			([xshift=-1pt]tw.north west) -- ([xshift=-1pt]tw.south west);
	\end{tikzpicture}%
	\end{center}%
}


\usepackage{contour}
\usepackage{ulem}

\renewcommand{\ULdepth}{1.8pt}
\contourlength{0.8pt}

\newcommand{\myuline}[1]{%
  \uline{\phantom{#1}}%
  \llap{\contour{white}{#1}}%
}

\usepackage[labelfont=bf,justification=centering]{caption}
\usepackage[labelfont=bf, justification=justified, labelformat=simple]{subcaption}
\captionsetup{position=top}
\captionsetup[subfigure]{justification=centering}
\renewcommand\thesubfigure{(\alph{subfigure})}

% MICRO-TYPOGRAPHY =========================================
% Loaded last so it applies to whichever font scheme was selected.
% Protrusion pulls punctuation slightly into the margin so ragged-right
% text edges look optically straight; expansion subtly varies glyph widths
% to even out word spacing. Both matter more on slides than in papers,
% because short ragged lines make spacing irregularities obvious.
\ifnum\slidemicrotype=1
	\usepackage[
		final,
		protrusion = true,
		expansion  = true,
		tracking   = true,
		kerning    = true,
		factor     = 1100,   % protrude a touch harder than the default
		stretch    = 12,
		shrink     = 12
	]{microtype}
	% Letterspace all-caps and small-caps runs, which otherwise look cramped.
	\SetTracking{encoding = {*}, shape = sc}{40}
	% Never letterspace or expand verbatim-like material.
	\microtypesetup{verbose = silent}
\fi

\usepackage{xkeyval}
\usepackage{todonotes}
\presetkeys{todonotes}{inline}{}